\section{Scope and Goals}
\label{sec:scope-goals}

\subsection{Project Goals}

The primary objective of this research is to design and implement an intelligent agent system that provides a natural language interface for Slurm cluster management. This system aims to democratize access to HPC resources by enabling users to interact with Slurm using conversational queries rather than complex command-line syntax. The following goals guide the development of this solution:

\textbf{Goal 1: Natural Language Interface for Slurm Operations.} The system shall enable users to perform common Slurm operations through natural language queries. Users should be able to check job status, view cluster resources, submit jobs, and monitor queues without explicit knowledge of Slurm command syntax. The agent must interpret user intent accurately and translate queries into appropriate Slurm commands.

\textbf{Goal 2: Multi-Agent Architecture with Specialized Capabilities.} The system shall implement a multi-agent architecture where specialized agents handle different aspects of cluster management. This includes agents for command execution, data visualization, and web search. The handoff mechanism between agents should be seamless and transparent to the user.

\textbf{Goal 3: Standardized Tool Integration via MCP.} The system shall utilize the Model Context Protocol to provide standardized tool interfaces between the LLM agents and Slurm infrastructure. The MCP server should expose Slurm functionality as discrete tools that can be invoked by the agent, ensuring clean separation between the AI layer and system operations.

\textbf{Goal 4: Visualization of Cluster Metrics.} The system shall provide visualization capabilities for cluster data, enabling users to request charts and graphs showing resource utilization, job distributions, and other metrics. Visualizations should be generated dynamically based on user queries and current cluster state.

\textbf{Goal 5: Real-Time Streaming Responses.} The system shall support real-time streaming of agent responses using Server-Sent Events (SSE), providing users with immediate feedback during long-running operations. This improves user experience by eliminating the perception of system unresponsiveness.

\textbf{Goal 6: Web-Based User Interface.} The system shall provide an intuitive web-based interface that supports markdown rendering, code highlighting, and inline visualization display. The interface should be accessible and usable without specialized client software.

\subsection{Project Scope}

The scope of this project encompasses the design, implementation, and evaluation of the Slurm Agent system. The following boundaries define what is included and excluded from the project:

\textbf{In Scope:}
\begin{itemize}
    \item Development of an MCP server exposing Slurm functionality as tools
    \item Implementation of a multi-agent system using the OpenAI Agents SDK
    \item Design and implementation of a web-based frontend interface
    \item Integration of chart visualization capabilities using matplotlib
    \item Implementation of web search functionality for documentation assistance
    \item Real-time response streaming using SSE
    \item Testing on a single-node Slurm cluster configuration
\end{itemize}

% \textbf{Out of Scope:}
% \begin{itemize}
%     \item Development of new Slurm features or modifications to Slurm core
%     \item Multi-user authentication and authorization systems
%     \item Production deployment considerations and security hardening
%     \item Performance benchmarking on large-scale production clusters
%     \item Training or fine-tuning of custom language models
%     \item Mobile application development
% \end{itemize}

This project focuses on demonstrating the feasibility and effectiveness of using LLM agents with MCP for HPC cluster management, providing a foundation for future research and development in this domain.
